\chapter*{Abstract}
\addcontentsline{toc}{chapter}{Abstract}
\markboth{Abstract}{Abstract}

This dissertation investigates capitalist reproduction when capital accumulation, productive capacity, and macroeconomic stability cannot be assumed to move proportionally or independently of distribution, institutions, and external constraints. Productive capacity is not directly observed, yet its identification determines the denominator of capacity utilization. Three essays connect this measurement problem to the structural conditions of capacity formation and to the historical operation of monetary and external constraints in the periphery.

The first essay critically replicates an accounting-based capacity measure for the US corporate sector over 1947--2011. The full-sample bivariate output--capital system does not sustain the required long-run properties. A relation survives when the rate of exploitation and historical breaks enter the system, with retained transformation elasticities below unity. This result challenges the treatment of productive capacity as an autonomous technical benchmark.

The second essay develops the principal structural mechanism connecting accumulation to capacity through distribution, choice of technique, internally divided firms, and capital composition. Long-run estimates for the United States over 1931--2024 and Chile over 1943--2010 indicate non-unitary transformation. In Chile, the estimated wage-share elasticity of the capital composition gap falls from approximately 3.29 to 0.14 when the external constraint binds. The comparison is analytically asymmetric: import dependence and foreign-exchange availability alter the feasible transformation process in the periphery.

The third essay examines the crisis of Chile's Unidad Popular through dependency theory, historical reconstruction, an open-economy central-bank balance-sheet framework, and monthly dynamic evidence. It distinguishes global monetary change, commodity-price movements, Chile-specific financial pressure, domestic conflict, and central-bank agency. In the adverse solvency-growth state, reverse monetary accommodation is more persistent and cumulatively larger than forward transmission, which remains significant. These findings are consistent with endogenous accommodation under external stress. They establish directional precedence and conditional responses, not experimental counterfactual causal identification.

Together, the essays connect capacity measurement to the relations governing its formation and show why peripheral reproduction also requires attention to international money and institutional choice. Their different designs do not statistically validate one another or establish a universal crisis law. They provide complementary explanations of how accumulation and macroeconomic reproduction depend on historically specific distributive, productive, and international relations.
